%Copyright 2021 Jean-Michel Bruel
%This program is free template, under the terms of the attached License of this repo. 
%Based on https://github.com/jpeisenbarth/SRS-Tex
%Itself based on the code of Yiannis Lazarides
%http://tex.stackexchange.com/questions/42602/software-requirements-specification-with-latex
%http://tex.stackexchange.com/users/963/yiannis-lazarides
%Also based on the template of Karl E. Wiegers
%http://www.se.rit.edu/~emad/teaching/slides/srs_template_sep14.pdf
%http://karlwiegers.com
\documentclass{scrreprt}
\usepackage{method}		% Specific command from Handbook
\usepackage{minitoc}	% For table of content in chapters
\usepackage{listings}
\usepackage{underscore}
\usepackage[bookmarks=true]{hyperref}
\usepackage[utf8]{inputenc}
\usepackage[english]{babel}
\usepackage{float}
\hypersetup{
    bookmarks=false,    % show bookmarks bar?
    pdftitle={Software Requirement Specification},    % title
    pdfauthor={Jean-Philippe Eisenbarth},                     % author
    pdfsubject={TeX and LaTeX},                        % subject of the document
    pdfkeywords={TeX, LaTeX, graphics, images}, % list of keywords
    colorlinks=true,       % false: boxed links; true: colored links
    linkcolor=blue,       % color of internal links
    citecolor=black,       % color of links to bibliography
    filecolor=black,        % color of file links
    urlcolor=purple,        % color of external links
    linktoc=page            % only page is linked
}%
\def\myversion{1.24.8 }
\date{}
%\title{%

%}
\usepackage{hyperref}
\newcommand*{\nsection}[1]{
    \section*{#1}
    \addcontentsline{toc}{section}{#1}
}

%--------------------- Numbering ---------------
\usepackage{amsthm}
\usepackage{xassoccnt}
%\newtheorem{req}{Requirement}[section]        
\newtheorem{req}{Requirement}        
\theoremstyle{definition}        
\newtheorem{constraint}{Constraint}
\newtheorem{goal}{Goal}
\DeclareCoupledCountersGroup{theorems}
\DeclareCoupledCounters[name=theorems]{req,constraint,goal}
\setcounter{goal}{0}
%--------------------- Comments ---------------
\newcommand{\mynote}[3][black]{\textcolor{#1}{\fbox{\bfseries\sffamily\scriptsize{#2}}
{\small\textsf{\emph{#3}}}}}
\newcommand{\comment}[1]{\mynote[red]{Comment:}{#1}}
%\newcommand{\comment}[1]{}
\newcommand{\notEmpty}{\comment{This chapter should not be empty!}}

\begin{document}

\begin{flushright}
    \rule{16cm}{5pt}\vskip1cm
    \begin{bfseries}
        \Huge{SOFTWARE REQUIREMENTS\\ SPECIFICATION}\\
        \vspace{1.9cm}
        for\\
        \vspace{1.9cm}
        $<$Project$>$\\
        \vspace{1.9cm}
        \LARGE{Version \myversion approved}\\
        \vspace{1.9cm}
        Prepared by $<$author$>$\\
        \vspace{1.9cm}
        $<$Organization$>$\\
        \vspace{1.9cm}
        \today\\
    \end{bfseries}
\end{flushright}

%===============================================================
\chapter*{Revision History}
%===============================================================

\begin{center}
    \begin{tabular}{|c|c|c|c|}
        \hline
	    Name        & Date          & Reason For Changes & Version\\
        \hline
        \hline
	    J.-M. Bruel & 2021-01-22    & First Draft & 1.0\\
        \hline
	    J.-M. Bruel & 2023-01-28    & Check after publication of the Handbook & 1.23\\
        \hline
	    J.-M. Bruel & 2023-06-12    & Add reqs automated numbering & 1.23.1 \\
        \hline
	    J.-M. Bruel & 2023-08-25    & Add Minimum Requirements Outcome Principle & 1.23.8 \\
        \hline
	    J.-M. Bruel & 2023-12-22    & Remove section numbers & 1.23.12 \\
        \hline
	    J.-M. Bruel & 2024-08-01    & Add warning about non empty chapters  & \myversion \\
        \hline
    \end{tabular}
\end{center}

This document follows the requirements documentation structure presented in the \href{ https://link.springer.com/content/pdf/10.1007/978-3-031-06739-6.pdf}{Handbook of requirements and business analysis}, by Bertrand Meyer.

%===============================================================
\dominitoc% Initialization
\tableofcontents
\adjustmtc
%===============================================================

%===============================================================
\addcontentsline{toc}{chapter}{Goals book}
\chapter*{Goals}

\minitoc% Creating an actual minitoc
%===============================================================

\comment{Goals are "needs of the target organization, which the system will address". While the development team is the principal user of the other books, the Goals book addresses a wider audience: essentially, all stakeholders.}

\input{Goals/G1.txt}
\input{Goals/G2.txt}
\input{Goals/G3.txt}
\input{Goals/G4.txt}
\input{Goals/G5.txt}
\input{Goals/G6.txt}
\input{Goals/G7.txt}

%===============================================================
\chapter*{Environment}
\addcontentsline{toc}{chapter}{Environment book}
\minitoc% Creating an actual minitoc
%===============================================================

\comment{The Environment book describes the application domain and external context, physical or virtual (or a mix), in which the system will operate.}

\input{Environment/E1.txt}
\input{Environment/E2.txt}
\input{Environment/E3.txt}
\input{Environment/E4.txt}
\input{Environment/E5.txt}
\input{Environment/E6.txt}

%===============================================================
\chapter*{System}
\addcontentsline{toc}{chapter}{System book}
\minitoc% Creating an actual minitoc
%===============================================================

\comment{The System book refines the Goal one by focusing on more detailed requirements about the system under development, mainly its constituents, behaviors and properties.}

\input{System/S1.txt}
\input{System/S2.txt}
\input{System/S3.txt}
\input{System/S4.txt}
\input{System/S5.txt}
\input{System/S6.txt}

%===============================================================
\chapter*{Project}
\addcontentsline{toc}{chapter}{Project book}
\minitoc% Creating an actual minitoc
%===============================================================

\comment{The Project book describes all the constraints and expectations not about the system itself, but about how to develop and produce it.}

\input{Project/P1.txt}
\input{Project/P2.txt}
\input{Project/P3.txt}
\input{Project/P4.txt}
\input{Project/P5.txt}
\input{Project/P6.txt}
\input{Project/P7.txt}


%===============================================================
\end{document}
%===============================================================
